% Source coverage: physical PDF pp. 215--234 (printed pp. 192--211), ending on the partial opening of physical p. 234 immediately before the Section 19.6 heading.
% Numbered equations: (19.5.1)--(19.5.71). Figure: 19.3. Citations: [12], [19], [21]--[30], [30a]. Footnotes: 2.
% Uncertainty: none.

\subsection{Effective Field Theories: Pions and Nucleons}
\label{sec:19.5}

In Section 19.2 we learned how to calculate the amplitude for emission of a single low energy Goldstone boson $B$ in a transition $\alpha\to\beta+B$ by applying the condition of current conservation to the matrix element of the symmetry current between the states $\alpha$ and $\beta$. In this calculation we never had to use any information about the details of the broken symmetry group; current conservation was all we needed. A new element enters if we wish to calculate the matrix element for the emission and/or absorption of \emph{two} Goldstone bosons, as for instance in a Goldstone boson scattering process. Here we must apply the condition of current conservation to a matrix element like
\[
\bra{\beta}T\{J_1^{\mu_1}(x_1),J_2^{\mu_2}(x_2)\}\ket{\alpha},
\]
where the states $\alpha$ and $\beta$ contain the other particles besides the two Goldstone bosons participating in the reaction. But when we let the divergence operator $\partial/\partial x_1^{\mu_1}$ act on this matrix element we encounter a non-zero contribution from the derivative of the functions $\theta(x_1^0-x_2^0)$ and $\theta(x_2^0-x_1^0)$ that appear in the definition of the time-ordered product $T\{\cdots\}$. This contribution is equal to the matrix element of an equal-time commutator $\delta(x_1^0-x_2^0)[J_1^0(x_1),J_2^{\mu_2}(x_2)]$, whose value depends on the commutation relations of the group algebra. This makes such multi-Goldstone-boson processes specially interesting, because they can be used to decide experimentally on the nature of the broken symmetry group, in a way that is not possible for processes involving just a single Goldstone boson. Because of the appearance of such current commutators, this approach is known as the method of current algebra~\hyperref[ch19-ref-21]{[21]}.

The current algebra method was used in early calculations of multi-Goldstone-boson amplitudes~\hyperref[ch19-ref-22]{[22]}. Unfortunately, such calculations are tedious, especially when three or more Goldstone bosons are involved, and it was also difficult to see how to deal with symmetries like the chiral symmetry of quantum chromodynamics that are not exact. For this reason a different and more physical calculational technique was introduced~\hyperref[ch19-ref-23]{[23]}, based on the use of effective Lagrangians: We simply calculate the Goldstone boson amplitudes by the methods of perturbation theory, using some Lagrangian for the Goldstone bosons and the other particles in the states $\alpha$ and $\beta$ that obeys the assumed broken symmetry.

Originally the justification for the effective Lagrangian procedure was based on current algebra. By using current algebra one could see that the amplitude for emission of a set of low energy Goldstone bosons in a process $\alpha\to\beta+B_1+B_2+\cdots$ is fixed once one knows the equal-time commutation relations of the currents associated with the broken symmetries as well as the matrix element for the process $\alpha\to\beta$ and the matrix elements of the currents between various one-particle states. We know that a Lagrangian that respects the broken symmetry will allow the construction of conserved currents with the appropriate equal-time commutators by Noether's technique, so if we simply calculate the low energy Goldstone boson amplitudes with such a Lagrangian and insert the correct values for $M_{\beta\alpha}$ and the one-particle matrix elements of the currents, we must get the same answer as provided by current algebra. In the case of interactions among Goldstone bosons alone, the states $\alpha$ and $\beta$ may both be taken as the vacuum, and we do not need any extra information beyond the matrix element $F$ between a Goldstone boson state and the vacuum.

In the first example of this sort~\hyperref[ch19-ref-23]{[23]}, the starting point was the Lagrangian of the $\sigma$-model~\hyperref[ch19-ref-24]{[24]}. Restricting our attention for the moment to the bosonic sector of this model, its Lagrangian is the $SO(4)$-invariant one used as an example in Section 19.2:
\begin{equation}
\mathcal L
=-\frac{1}{2}\partial_\mu\phi_n\partial^\mu\phi_n
-\frac{\mathcal M^2}{2}\phi_n\phi_n
-\frac{\lambda}{4}(\phi_n\phi_n)^2,
\tag{19.5.1}\label{eq:19.5.1}
\end{equation}
where $n$ is understood to be summed over the values $1,2,3,4$, with $\boldsymbol{\phi}$ an isovector pseudoscalar field and $\phi_4$ an isoscalar scalar field.

The immediate problem faced by any sort of effective Lagrangian is that in order to use it to calculate scattering amplitudes, we must either include all Feynman diagrams of all orders of perturbation theory, or else find some rationale for dropping higher-order diagrams. We can find no such rationale with Lagrangians like \eqref{eq:19.5.1}, in which the broken symmetry is realized through linear transformations on the various fields. Fortunately any such Lagrangian can be recast in a form that allows the use of Feynman diagrams to generate an expansion for scattering amplitudes in powers of Goldstone boson energies. To do this, we perform a symmetry transformation at every point in spacetime that eliminates the fields corresponding to the Goldstone bosons of the theory. The Goldstone boson degrees of freedom reappear in the transformed theory as the parameters of this symmetry transformation. However, since the Lagrangian is invariant under spacetime-independent symmetry transformations, it cannot have any dependence on the new Goldstone boson fields when they are constant, and so every term in the Lagrangian that involves these new Goldstone boson fields must contain at least one spacetime derivative of the field. These derivatives introduce factors of the Goldstone boson energy when we calculate S-matrix elements for Goldstone boson reactions, and as we shall see we can use the Lagrangian to construct a series for the S-matrix elements in powers of these energies.

For instance, to recast the Lagrangian \eqref{eq:19.5.1} in a useful form, we write the four-vector $\phi_n$ as a chiral rotation $R$ acting on a four-vector $\{0,0,0,\sigma\}$ whose first three components (the Goldstone part of $\phi_n$) vanish:
\begin{equation}
\phi_n(x)=R_{n4}(x)\sigma(x)
\tag{19.5.2}\label{eq:19.5.2}
\end{equation}
with $R_{nm}(x)$ an orthogonal matrix
\begin{equation}
R^T(x)R(x)=\mathbf{1}
\tag{19.5.3}\label{eq:19.5.3}
\end{equation}
and therefore
\begin{equation}
\sigma(x)=\sqrt{\sum_n\phi_n(x)^2}.
\tag{19.5.4}\label{eq:19.5.4}
\end{equation}
The Lagrangian \eqref{eq:19.5.1} then becomes
\begin{equation}
\mathcal L
=-\frac{1}{2}\sum_{n=1}^4
\left(R_{n4}\partial_\mu\sigma+\sigma\partial_\mu R_{n4}\right)^2
-\frac{1}{2}\mathcal M^2\sigma^2-\frac{\lambda}{4}\sigma^4.
\tag{19.5.5}\label{eq:19.5.5}
\end{equation}
Because $R$ is an orthogonal matrix, the $\partial_\mu\sigma\partial^\mu\sigma$ term is $R$-independent, and the cross-term vanishes
\[
\sum_nR_{n4}^2=1,\qquad
\sum_nR_{n4}\partial_\mu R_{n4}
=\frac{1}{2}\partial_\mu\sum_nR_{n4}^2=0,
\]
so $\mathcal L$ becomes
\begin{equation}
\mathcal L
=-\frac{1}{2}\partial_\mu\sigma\partial^\mu\sigma
-\frac{1}{2}\sigma^2\sum_{n=1}^4
\partial^\mu R_{n4}\partial_\mu R_{n4}
-\frac{1}{2}\mathcal M^2\sigma^2-\frac{\lambda}{4}\sigma^4.
\tag{19.5.6}\label{eq:19.5.6}
\end{equation}
If $\mathcal M^2$ is negative then $\sigma$ has a non-vanishing vacuum expectation value, given in lowest order by the position of the minimum of the sum of the last two terms, at $\sigma=|\mathcal M|/\sqrt{\lambda}$.

In place of the field variables $\phi_n$, our variables now are $\sigma-|\mathcal M|/\sqrt{\lambda}$ and whatever other variables are needed to parameterize the rotation $R$. For instance, these parameters could be simply chosen as the $R_{a4}$ themselves (where $a,b,\ldots$ from now are isovector indices running over the values $1,2,3$), with $R_{44}$ given by the condition that $R$ is orthogonal. A different parameterization will give simpler final results, and was historically the first to be used for these purposes. It is to define
\begin{equation}
\zeta_a=\frac{\phi_a}{\phi_4+\sigma}
\tag{19.5.7}\label{eq:19.5.7}
\end{equation}
and take
\begin{equation}
R_{a4}=\frac{2\zeta_a}{1+\boldsymbol{\zeta}^2}=-R_{4a},
\qquad
R_{44}=\frac{1-\boldsymbol{\zeta}^2}{1+\boldsymbol{\zeta}^2},
\qquad
R_{ab}=\delta_{ab}-\frac{2\zeta_a\zeta_b}{1+\boldsymbol{\zeta}^2},
\tag{19.5.8}\label{eq:19.5.8}
\end{equation}
so that
\begin{equation}
\phi_a/\sigma=R_{a4}=\frac{2\zeta_a}{1+\boldsymbol{\zeta}^2},
\qquad
\phi_4/\sigma=R_{44}=\frac{1-\boldsymbol{\zeta}^2}{1+\boldsymbol{\zeta}^2}.
\tag{19.5.9}\label{eq:19.5.9}
\end{equation}

Then the Lagrangian \eqref{eq:19.5.6} becomes
\begin{equation}
\mathcal L
=-\frac{1}{2}\partial_\mu\sigma\partial^\mu\sigma
-2\sigma^2\boldsymbol D_\mu\cdot\boldsymbol D^\mu
-\frac{1}{2}\mathcal M^2\sigma^2-\frac{\lambda}{4}\sigma^4,
\tag{19.5.10}\label{eq:19.5.10}
\end{equation}
where
\begin{equation}
\boldsymbol D_\mu\equiv
\frac{\partial_\mu\boldsymbol{\zeta}}{1+\boldsymbol{\zeta}^2}.
\tag{19.5.11}\label{eq:19.5.11}
\end{equation}
Whatever parameterization we use, it is clear that the fields $\boldsymbol{\zeta}$ describe particles of zero mass, whose interactions all involve field derivatives. These are (up to a normalization) our new pion fields.

Despite appearances, this Lagrangian is still invariant under $SO(4)$, but with $SO(4)$ now realized non-linearly. Under an isospin transformation with infinitesimal parameters $\boldsymbol{\theta}$, the field $\boldsymbol{\zeta}$ simply rotates like an ordinary isovector, and $\sigma$ is an isoscalar:
\begin{equation}
\delta\boldsymbol{\zeta}
=\boldsymbol{\theta}\cdot\boldsymbol{\zeta},
\qquad
\delta\sigma=0,
\tag{19.5.12}\label{eq:19.5.12}
\end{equation}
so the Lagrangian \eqref{eq:19.5.10} is manifestly isospin-invariant. On the other hand, under the broken symmetry transformations parameterized by an infinitesimal vector $\boldsymbol{\epsilon}$, the original fields transform according to
\begin{equation}
\delta\boldsymbol{\phi}=2\boldsymbol{\epsilon}\phi_4,
\qquad
\delta\phi_4=-2\boldsymbol{\epsilon}\cdot\boldsymbol{\phi}.
\tag{19.5.13}\label{eq:19.5.13}
\end{equation}
From \eqref{eq:19.5.7}, we then find
\begin{equation}
\delta\boldsymbol{\zeta}
=\boldsymbol{\epsilon}(1-\boldsymbol{\zeta}^2)
+2\boldsymbol{\zeta}(\boldsymbol{\epsilon}\cdot\boldsymbol{\zeta}),
\qquad
\delta\sigma=0.
\tag{19.5.14}\label{eq:19.5.14}
\end{equation}
The Lagrangian \eqref{eq:19.5.10} is invariant under the broken symmetry transformation \eqref{eq:19.5.14} because $\boldsymbol D_\mu$ undergoes a linear (though field-dependent) isospin rotation~\hyperref[ch19-ref-25]{[25]}:
\begin{equation}
\delta\boldsymbol D_\mu
=2(\boldsymbol{\zeta}\cdot\boldsymbol{\epsilon})
\cdot\boldsymbol D_\mu
\tag{19.5.15}\label{eq:19.5.15}
\end{equation}
and \eqref{eq:19.5.10} is isospin-invariant. Because of the transformation rule \eqref{eq:19.5.15}, $\boldsymbol D_\mu$ is often called the \emph{covariant derivative} of the pion field.

The transformation rules \eqref{eq:19.5.12} and \eqref{eq:19.5.14} specify what is called a \emph{non-linear realization} of the group $SU(2)\cdot SU(2)$~\hyperref[ch19-ref-25]{[25]}. The general theory of non-linear realizations of Lie groups is given in the next section; we shall show there that, up to field redefinitions, the transformation rules \eqref{eq:19.5.12} and \eqref{eq:19.5.14} provide the most general realization of $SU(2)\cdot SU(2)$ in which the isospin $SU(2)$ subgroup is realized linearly on $\boldsymbol{\zeta}$.

We see that each interaction of these new pion fields is accompanied with a spacetime derivative, so that the effective coupling is weak for low pion energies. (This remark will be made more precise below.) Therefore for sufficiently small pion energies we may use this Lagrangian in the tree approximation to reproduce the soft-pion theorems of current algebra. For this purpose, it is only necessary that the Lagrangian be $SO(4)$-invariant. But since the $\sigma$ field is an $SO(4)$ scalar, it plays no role in maintaining the $SO(4)$ invariance of the Lagrangian, and may be simply discarded\footnote{Alternatively, we can pass to the limit where $\mathcal M$ and $\lambda$ go to infinity together, keeping the expectation value of $\sigma$ constant.}. Of course, this procedure changes the physical content of the theory, but it does not change the amplitudes given by soft-pion theorems. The Lagrangian \eqref{eq:19.5.10} then simplifies to
\begin{equation}
\mathcal L
=-\frac{F^2}{2}\boldsymbol D_\mu\cdot\boldsymbol D^\mu
=-\frac{F^2}{2}
\frac{\partial_\mu\boldsymbol{\zeta}\cdot
\partial^\mu\boldsymbol{\zeta}}
{(1+\boldsymbol{\zeta}^2)^2},
\tag{19.5.16}\label{eq:19.5.16}
\end{equation}
where $F=2\langle\sigma\rangle=2|\mathcal M|/\sqrt{\lambda}$. (As we shall soon see, this $F$ is the same as the constant $F_\pi$ discussed in the previous section.) For many purposes it is more convenient to work with a conventionally normalized pion field
\begin{equation}
\boldsymbol{\pi}\equiv F\boldsymbol{\zeta}
\tag{19.5.17}\label{eq:19.5.17}
\end{equation}
for which the Lagrangian \eqref{eq:19.5.16} reads
\begin{equation}
\mathcal L
=-\frac{1}{2}
\frac{\partial_\mu\boldsymbol{\pi}\cdot
\partial^\mu\boldsymbol{\pi}}
{(1+\boldsymbol{\pi}^2/F^2)^2}.
\tag{19.5.18}\label{eq:19.5.18}
\end{equation}

The factor $1/F$ acts as a coupling parameter that accompanies the interaction of each additional pion. Eq.~\eqref{eq:19.5.18} describes what is often called the `non-linear $\sigma$-model,' for the special case of $SU(2)\cdot SU(2)$ spontaneously broken to $SU(2)$.

An important point: to derive \eqref{eq:19.5.18} it was not really necessary to start with the `linear $\sigma$-model' Lagrangian \eqref{eq:19.5.1}. Indeed, we did not need to start with \emph{any} specific theory. Eq.~\eqref{eq:19.5.18} can be used simply because it is invariant under the $SO(4)$ transformations \eqref{eq:19.5.12}, \eqref{eq:19.5.14}, and current algebra tells us that this is all we need to get the right results for low energy pion reaction amplitudes.

Some years after the introduction of effective Lagrangians for soft pions, there emerged a different justification for the effective field theory technique~\hyperref[ch19-ref-26]{[26]}, one that does not rely on current algebra and allows calculations that are not limited to the limit of vanishingly small Goldstone boson energies. It is based on the realization (not yet formally embodied in a theorem) that when we calculate a physical amplitude from Feynman diagrams using the most general Lagrangian that involves the relevant degrees of freedom and satisfies the assumed symmetries of the theory, we are simply constructing the most general amplitude that is consistent with general principles of relativity, quantum mechanics, and the assumed symmetries. This was the point of view underlying Volume I of this book. In the present context, the `relevant' degrees of freedom are the Goldstone bosons themselves, together with the particles in the states $\alpha$ and $\beta$ and any other particle states that can be produced from them by interactions with low energy Goldstone bosons. By invoking this justification for effective field theories, we are freed from any need to wrestle with the complications of current algebra. More importantly, the modern effective field theory approach yields results that take us beyond the extreme low energy limit and allows a systematic treatment of any intrinsic symmetry breaking.

According to this approach, in order to calculate pion interaction amplitudes to any desired order in pion energies, we must use the most general Lagrangian involving a pion field $\boldsymbol{\zeta}$ that transforms according to the rules \eqref{eq:19.5.12} and \eqref{eq:19.5.14}:
\begin{equation}
\mathcal L_{\mathrm{eff}}
=-\frac{F^2}{2}\boldsymbol D_\mu\cdot\boldsymbol D^\mu
-\frac{c_4}{4}(\boldsymbol D_\mu\cdot\boldsymbol D^\mu)^2
-\frac{c'_4}{4}
(\boldsymbol D_\mu\cdot\boldsymbol D_\nu)
(\boldsymbol D^\mu\cdot\boldsymbol D^\nu)-\cdots.
\tag{19.5.19}\label{eq:19.5.19}
\end{equation}
The terms indicated by $\cdots$ will contain higher powers of the covariant derivative $\boldsymbol D_\mu$, or higher covariant derivatives, whose general structure is described in the next section. The coefficients $c_4$ and $c'_4$ are dimensionless, and all higher terms have coefficients with the dimensionality of negative powers of mass.

Consider a general process involving arbitrary numbers of incoming and outgoing pions. We suppose that their energies and momenta are all at most of some order $Q$, which is small compared with a typical quantum chromodynamics energy scale (say, the nucleon or $\rho$ mass). Even though Lagrangians like \eqref{eq:19.5.19} are not renormalizable in the usual sense, we saw in Section 12.3 that such Lagrangians can yield finite results as long as they contain all possible terms allowed by symmetries, for then there will be a counterterm available to cancel every infinity. If we define the renormalized values of the constants $F^2,c_4,c'_4,\ldots$ by specifying the values of various Goldstone boson scattering amplitudes at energies of order $Q$, then the integrals in momentum-space Feynman diagrams will be dominated by contributions from virtual momenta which are also of order $Q$ (because renormalization makes them finite, and there is no other possible effective cut-off in the theory.) We can then develop perturbation theory as a power series expansion in $Q$.

Each derivative and hence each $\boldsymbol D_\mu$ in each interaction vertex contributes one factor of $Q$ to the order of magnitude of the diagram; each internal pion propagator contributes a factor $Q^{-2}$; and each integration volume $d^4q$ associated with the loops of the diagram contributes a factor $Q^4$; so a general connected diagram makes a contribution of order $Q^\nu$, where
\begin{equation}
\nu=\sum_iV_i d_i-2I+4L.
\tag{19.5.20}\label{eq:19.5.20}
\end{equation}
Here $d_i$ is the number of derivatives in an interaction of type $i$, $V_i$ is the number of interaction vertices of type $i$ in the diagram, $I$ is the number of internal pion lines, and $L$ is the number of loops. These quantities are related by the familiar topological identity (see Eq.~(4.4.7)):
\begin{equation}
L=I-\sum_iV_i+1,
\tag{19.5.21}\label{eq:19.5.21}
\end{equation}
so we can eliminate $I$ and write
\begin{equation}
\nu=\sum_iV_i(d_i-2)+2L+2.
\tag{19.5.22}\label{eq:19.5.22}
\end{equation}

The point of this is that each term here is positive; every interaction in \eqref{eq:19.5.19} has at least two derivatives, and of course $L\geq0$. Therefore the leading term of each process is of order $Q^2$, and arises solely from \emph{tree} graphs (that is, $L=0$) constructed solely from the term in \eqref{eq:19.5.19} with only two derivatives (that is, $V_i=0$ for $d_i\neq2$):
\begin{equation}
\mathcal L_2
=-\frac{F^2}{2}\boldsymbol D_\mu\cdot\boldsymbol D^\mu
=-\frac{1}{2}
\frac{\partial_\mu\boldsymbol{\pi}\cdot
\partial^\mu\boldsymbol{\pi}}
{(1+\boldsymbol{\pi}^2/F^2)^2}.
\tag{19.5.23}\label{eq:19.5.23}
\end{equation}
For instance, the invariant amplitude $M$ that appears in the pion-pion scattering S-matrix element
\begin{equation}
S=i(2\pi)^4\delta^4(p_A+p_B-p_C-p_D)M
(2\pi)^{-6}(16E_AE_BE_CE_D)^{-1/2}
\tag{19.5.24}\label{eq:19.5.24}
\end{equation}
is given to this order by
\begin{align}
M_{abcd}^{(\nu=2)}
=4F^{-2}\Big[&
\delta_{ab}\delta_{cd}(-p_A\cdot p_B-p_C\cdot p_D)
\notag\\
&+\delta_{ac}\delta_{bd}(p_A\cdot p_C+p_B\cdot p_D)
\notag\\
&+\delta_{ad}\delta_{bc}(p_A\cdot p_D+p_B\cdot p_C)
\Big].
\tag{19.5.25}\label{eq:19.5.25}
\end{align}
Here $a,b,c,d$ are the isovector indices of the pions $A,B,C,D$, respectively. (The effect of the finite pion mass will be taken up a little later in this section.) Using \eqref{eq:19.5.25} as the leading term does not depend on any assumption of the smallness of the coupling constant $\lambda$ in the original Lagrangian \eqref{eq:19.5.1}, or even on the validity of this formula for the Lagrangian, but only on the assumed smallness of the typical pion energy $Q$.

The next term in the amplitude for any Goldstone boson reaction will be of order $Q^4$, and arises both from one-loop graphs involving only the Lagrangian \eqref{eq:19.5.16}, and also from tree graphs constructed solely from the interaction \eqref{eq:19.5.16} plus a single vertex arising from the $d=4$ terms in \eqref{eq:19.5.19}:
\begin{align}
M_{abcd}^{(\nu=4)}
={}&\frac{\delta_{ab}\delta_{cd}}{F^4}
\Bigg[
-\frac{1}{2\pi^2}s^2\ln(-s)
-\frac{1}{12\pi^2}(u^2-s^2+3t^2)\ln(-t)
\notag\\
&\hspace{1.8cm}
-\frac{1}{12\pi^2}(t^2-s^2+3u^2)\ln(-u)
+\frac{1}{3\pi^2}(s^2+t^2+u^2)\ln\Lambda^2
\notag\\
&\hspace{1.8cm}
-\frac{1}{2}c_4s^2-\frac{1}{4}c'_4(t^2+u^2)
\Bigg]
+\text{crossed terms},
\tag{19.5.26}\label{eq:19.5.26}
\end{align}
where `crossed terms' denotes terms given by the interchanges $B\leftrightarrow C$ and $B\leftrightarrow D$, and $s,t$, and $u$ are the Mandelstam variables
\[
s=-(p_A+p_B)^2,\qquad
t=-(p_A-p_C)^2,\qquad
u=-(p_A-p_D)^2.
\]
The dependence of this result on the cut-off $\Lambda$ can be eliminated by a redefinition of the constants $c_4$ and $c'_4$. The renormalized couplings are
\begin{equation}
c_{4R}=c_4-\frac{2}{3\pi^2}\ln\left(\frac{\Lambda^2}{\mu^2}\right),
\tag{19.5.27}\label{eq:19.5.27}
\end{equation}
\begin{equation}
c'_{4R}=c'_4-\frac{4}{3\pi^2}\ln\left(\frac{\Lambda^2}{\mu^2}\right),
\tag{19.5.28}\label{eq:19.5.28}
\end{equation}
where $\mu$ is an arbitrary renormalization scale of order $Q$ inserted in order to make the logarithms well-defined. In terms of these renormalized couplings, the amplitude \eqref{eq:19.5.26} takes the form
\begin{align}
M_{abcd}^{(\nu=4)}
={}&\frac{\delta_{ab}\delta_{cd}}{F^4}
\Bigg[
-\frac{1}{2\pi^2}s^2\ln\left(\frac{-s}{\mu^2}\right)
\notag\\
&\hspace{0.8cm}
-\frac{1}{12\pi^2}(u^2-s^2+3t^2)
\ln\left(\frac{-t}{\mu^2}\right)
\notag\\
&\hspace{0.8cm}
-\frac{1}{12\pi^2}(t^2-s^2+3u^2)
\ln\left(\frac{-u}{\mu^2}\right)
-\frac{1}{2}c_{4R}s^2-\frac{1}{4}c'_{4R}(t^2+u^2)
\Bigg]
\notag\\
&+\text{crossed terms}.
\tag{19.5.29}\label{eq:19.5.29}
\end{align}
This sort of calculation can be carried to arbitrary orders in $Q$, always with the result that in each order we encounter a finite number of new couplings whose renormalization serves to eliminate the cut-off dependence of physical amplitudes. Notice that the ratio of the leading $\nu=2$ terms and the $\nu=4$ corrections is of order $Q^2/8\pi^2F^2$, so this is likely to be a useful expansion as long as the pion energies are all much less than an amount of order $2\pi F$.

This perturbative expansion may also be used to relate the ubiquitous parameter $F$ to the measured pion decay amplitude $F_\pi$. Recalling the transformation rule \eqref{eq:19.5.14} for $\boldsymbol{\zeta}$, the axial-vector current is given by
\[
\boldsymbol{\epsilon}\cdot\boldsymbol A^\mu
=-\frac{\partial\mathcal L}
{\partial(\partial_\mu\boldsymbol{\zeta})}
\cdot\delta\boldsymbol{\zeta}
\]
and so
\begin{equation}
\boldsymbol A^\mu
=-(1-\boldsymbol{\zeta}^2)
\frac{\partial\mathcal L}
{\partial(\partial_\mu\boldsymbol{\zeta})}
-2\boldsymbol{\zeta}\,
\boldsymbol{\zeta}\cdot
\left(
\frac{\partial\mathcal L}
{\partial(\partial_\mu\boldsymbol{\zeta})}
\right).
\tag{19.5.30}\label{eq:19.5.30}
\end{equation}
(This $\boldsymbol A^\mu$ is the Noether current of the symmetry generated by $2\boldsymbol{\epsilon}$, which on the nucleon doublet is represented by $2\gamma_5\boldsymbol t=\gamma_5\boldsymbol{\tau}$.) After integrating out all other fields, and using \eqref{eq:19.5.17} to express $\boldsymbol{\zeta}$ in terms of the canonically normalized pion field $\boldsymbol{\pi}$, we find
\begin{equation}
\boldsymbol A^\mu
=F\left[
\partial^\mu\boldsymbol{\pi}\,
\frac{1-\boldsymbol{\pi}^2/F^2}
{1+\boldsymbol{\pi}^2/F^2}
+\frac{2\boldsymbol{\pi}
(\boldsymbol{\pi}\cdot\partial^\mu\boldsymbol{\pi})}
{F^2(1+\boldsymbol{\pi}^2/F^2)^2}
\right]+\cdots.
\tag{19.5.31}\label{eq:19.5.31}
\end{equation}
Using our perturbative expansion in powers of the pion energy to calculate $\bra{\Omega}\boldsymbol A^\mu\ket{\boldsymbol{\pi}}$, we see that in lowest order the pion decay amplitude here is
\begin{equation}
F_\pi=F.
\tag{19.5.32}\label{eq:19.5.32}
\end{equation}
Furthermore Lorentz invariance and \eqref{eq:19.5.22} tell us that the higher-order corrections must be proportional to powers of $p_\pi^2/F_\pi^2$, which vanishes for a massless pion, so \eqref{eq:19.5.32} is actually exact in the limit $m_\pi\to0$.

We may therefore guess that our perturbation expansion will be useful for pion energies that are less than an amount of order $2\pi F_\pi=1200\ \mathrm{MeV}$.

In order to make contact with experiment, we must deal with the fact that the pion mass is not zero. On the mass shell it is not possible for the time component of a pion four-momentum to be less than $m_\pi$, so in counting powers of the typical pion energy and momentum $Q$, we must regard $m_\pi$ as being of order $Q$. But we saw in the previous section that $m_\pi^2$ is proportional to a linear combination of quark masses, so our formula \eqref{eq:19.5.22} for the order $Q^\nu$ of a given Feynman diagram should be rewritten to read
\begin{equation}
\nu=\sum_iV_i(d_i+2m_i-2)+2L+2,
\tag{19.5.33}\label{eq:19.5.33}
\end{equation}
where $m_i$ is the number of factors of quark masses in the interaction of type $i$.

The interactions involving quark masses may be distinguished by their transformation properties under $SU(2)\cdot SU(2)$, or equivalently, under $SO(4)$. Eq.~\eqref{eq:19.4.34} shows that the terms of first order in quark masses transform as a linear combination of two scalars, the fourth component of a chiral four-vector $\Phi_n^+$, with coefficient $m_u+m_d$, and the third component of a different chiral four-vector $\Phi_n^-$, with coefficient $m_u-m_d$. Thus the terms $\Delta\mathcal L^+$ and $\Delta\mathcal L^-$ in the effective Lagrangian that are of first order in $m_u+m_d$ and $m_u-m_d$ must have the chiral transformation properties of the fourth and third components of chiral four-vectors, respectively, and of course be Lorentz scalars. One obvious candidate for a chiral four-vector whose fourth component is a scalar is the $\phi_n$ with which we started in this section. According to \eqref{eq:19.5.9}, its fourth component is just $\sigma(1-\boldsymbol{\zeta}^2)/(1+\boldsymbol{\zeta}^2)$. The factor $\sigma$ may be dropped, as it is a chiral scalar, and therefore has no effect on the chiral transformation properties of this quantity. We can fix the normalization of this term by requiring that the coefficient of $\boldsymbol{\pi}^2=F_\pi^2\boldsymbol{\zeta}^2$ be $-m_\pi^2/2$, so that, apart from an additive constant,
\begin{equation}
\Delta\mathcal L^+
=-\frac{m_\pi^2F_\pi^2}{2}
\frac{\boldsymbol{\zeta}^2}{1+\boldsymbol{\zeta}^2}
=-\frac{m_\pi^2}{2}
\frac{\boldsymbol{\pi}^2}
{1+\boldsymbol{\pi}^2/F_\pi^2}.
\tag{19.5.34}\label{eq:19.5.34}
\end{equation}
We shall see in the next section that this is the unique scalar function of the pion field without derivatives that transforms as the fourth component of a chiral four-vector. On the other hand, there is \emph{no} scalar function of the pion field without derivatives that transforms as the third component of a chiral four-vector because such a function, if the third component of a chiral four-vector, would have to be odd in the pion field, and therefore pseudoscalar rather than scalar. Thus \eqref{eq:19.5.34} is the only interaction with $d_i=0$ and $m_i=1$. It is striking that the isospin violating difference in the $u$ and $d$ quark masses has not only no effect on the pion masses, as we saw in the previous section, but also has no effect on any non-derivative multipion interaction.

We are now in a position to do a realistic calculation of the leading $\nu=2$ terms in the pion-pion scattering amplitude. According to \eqref{eq:19.5.33}, these terms arise only from tree graphs constructed from the $d=2$, $m=0$ pion interaction in \eqref{eq:19.5.23} or from the $d=0$, $m=1$ pion interaction in \eqref{eq:19.5.34}. To this order, the invariant amplitude $M$ defined by \eqref{eq:19.5.24} is now
\begin{equation}
M_{abcd}^{(\nu=2)}
=4F_\pi^{-2}\left[
\delta_{ab}\delta_{cd}(s-m_\pi^2)
+\delta_{ac}\delta_{bd}(t-m_\pi^2)
+\delta_{ad}\delta_{bc}(u-m_\pi^2)
\right].
\tag{19.5.35}\label{eq:19.5.35}
\end{equation}
In particular, at threshold $s=4m_\pi^2$, $t=u=0$, so
\begin{align}
M_{abcd}^{(\nu=2)}(\text{threshold})
&=4m_\pi^2F_\pi^{-2}
[3\delta_{ab}\delta_{cd}
-\delta_{ac}\delta_{bd}-\delta_{ad}\delta_{bc}]
\notag\\
&=4m_\pi^2F_\pi^{-2}
\left[7M_{ab,cd}^{(0)}-2M_{ab,cd}^{(2)}\right],
\tag{19.5.36}\label{eq:19.5.36}
\end{align}
where $M^{(0)}$ and $M^{(2)}$ are the appropriate tensors representing two-pion states with isospin $T=0$ or $T=2$:
\begin{equation}
M_{ab,cd}^{(0)}=\frac{1}{3}\delta_{ab}\delta_{cd},
\tag{19.5.37}\label{eq:19.5.37}
\end{equation}
\begin{equation}
M_{ab,cd}^{(2)}
=\frac{1}{2}\left(
\delta_{ac}\delta_{bd}+\delta_{ad}\delta_{bc}
-\frac{2}{3}\delta_{ab}\delta_{cd}
\right),
\tag{19.5.38}\label{eq:19.5.38}
\end{equation}
normalized so that $\operatorname{Tr}M^{(T)}=2T+1$. This result is usually expressed in terms of the scattering lengths. According to Sections 3.6 and 3.7, the scattering length $a_T$ for two pions in a state of isospin $T$ is given~\hyperref[ch19-ref-27]{[27]} by $1/32\pi m_\pi$ times the coefficient of $M^{(T)}$ in \eqref{eq:19.5.36}:
\[
a_0=\frac{7m_\pi}{8\pi F_\pi^2}=0.16m_\pi^{-1},
\qquad
a_2=-\frac{m_\pi}{4\pi F_\pi^2}=-0.046m_\pi^{-1}.
\]

The pion scattering lengths are difficult to measure, but careful study of processes like $\pi+N\to\pi+\pi+N$ and $K\to\pi+\pi+e^-+\nu$ has given the results~\hyperref[ch19-ref-28]{[28]} $(0.26\mathbin{\pm}0.05)m_\pi^{-1}$ and $(-0.028\mathbin{\pm}0.012)m_\pi^{-1}$ for $a_0$ and $a_2$, respectively, which are consistent with these theoretical values. Corrections of higher order in $m_\pi/2\pi F_\pi$ seem to improve the agreement.

This formalism can be extended to describe the interactions of pions with nucleons or other particles. The easiest approach is to suppose that we add a term involving the nucleon doublet field $N$ to the Lagrangian with which we started:
\begin{equation}
\mathcal L_N
=-\overline N\left(
\sl{\partial}
+g[\phi_4+2i\boldsymbol t\cdot\boldsymbol{\phi}\gamma_5]
\right)N,
\tag{19.5.39}\label{eq:19.5.39}
\end{equation}
where $\boldsymbol t$ is the isospin matrix vector for isospin $1/2$ (that is, half the Pauli spinor $\boldsymbol{\tau}$.) This is invariant under chiral $SU(2)\cdot SU(2)$ transformations
\begin{equation}
\delta\boldsymbol{\phi}=2\boldsymbol{\epsilon}\phi_4,
\qquad
\delta\phi_4=-2\boldsymbol{\epsilon}\cdot\boldsymbol{\phi},
\tag{19.5.40}\label{eq:19.5.40}
\end{equation}
\begin{equation}
\delta N=-2i\gamma_5\boldsymbol{\epsilon}\cdot\boldsymbol t\,N.
\tag{19.5.41}\label{eq:19.5.41}
\end{equation}
Now in eliminating the non-derivative pion couplings we must express the nucleon field $N$ as the $SO(4)$ rotation $R$ in the representation \eqref{eq:19.5.41} acting upon a new nucleon field $\widetilde N$
\begin{equation}
N\equiv
\frac{1-2i\gamma_5\boldsymbol t\cdot\boldsymbol{\zeta}}
{\sqrt{1+\boldsymbol{\zeta}^2}}\widetilde N
\tag{19.5.42}\label{eq:19.5.42}
\end{equation}
with $\boldsymbol{\zeta}$ given again by \eqref{eq:19.5.7}. With this transformation the non-derivative term in \eqref{eq:19.5.39} now depends only on $\widetilde N$ and $\sigma$:
\[
\overline N[\phi_4+2i\boldsymbol t\cdot\boldsymbol{\phi}\gamma_5]N
=\sigma\overline{\widetilde N}\widetilde N.
\]
On the other hand, the derivative term involves derivatives of the matrix in \eqref{eq:19.5.41}, yielding a nucleonic Lagrangian:
\begin{equation}
\mathcal L_N
=-\overline{\widetilde N}
\left[
\sl{\partial}+g\sigma
+2i\frac{\boldsymbol t\cdot
(\boldsymbol{\zeta}\cdot\sl{\partial}\boldsymbol{\zeta})}
{1+\boldsymbol{\zeta}^2}
+2i\gamma_5\boldsymbol t\cdot\sl{\boldsymbol D}
\right]\widetilde N
\tag{19.5.43}\label{eq:19.5.43}
\end{equation}
or in terms of the canonically normalized pion field \eqref{eq:19.5.17}:
\begin{equation}
\mathcal L_N
=-\overline{\widetilde N}
\left[
\sl{\partial}+g\sigma
+\frac{2i\boldsymbol t\cdot
(\boldsymbol{\pi}\cdot\sl{\partial}\boldsymbol{\pi})}
{F_\pi^2[1+\boldsymbol{\pi}^2/F_\pi^2]}
+\frac{2i\gamma_5\boldsymbol t\cdot
\sl{\partial}\boldsymbol{\pi}}
{F_\pi[1+\boldsymbol{\pi}^2/F_\pi^2]}
\right]\widetilde N.
\tag{19.5.44}\label{eq:19.5.44}
\end{equation}

Since $\sigma$ has a non-vanishing vacuum expectation value, we see that the nucleon here has a non-zero rest mass, a result that would be prohibited if the symmetry under the transformations \eqref{eq:19.5.40}, \eqref{eq:19.5.41} were unbroken.

By its construction, the Lagrangian \eqref{eq:19.5.44} is chiral-invariant. This can also be seen directly, using the previously obtained chiral transformation properties of $\sigma$ and $\boldsymbol{\pi}$, and now also noting that under the chiral transformations \eqref{eq:19.5.40}, \eqref{eq:19.5.41}, the new nucleon field defined by \eqref{eq:19.5.42} transforms as
\begin{equation}
\delta\widetilde N
=2i\boldsymbol t\cdot
[\boldsymbol{\zeta}\cdot\boldsymbol{\epsilon}]
\widetilde N.
\tag{19.5.45}\label{eq:19.5.45}
\end{equation}
That is, under a chiral transformation $\widetilde N$ simply undergoes the same isospin rotation \eqref{eq:19.5.15} as the quantity $\boldsymbol D_\mu$, but of course in the $T=1/2$ representation. The isospin rotation parameter $\boldsymbol{\zeta}\cdot\boldsymbol{\epsilon}$ is spacetime-dependent, so derivatives of $\widetilde N$ do not have the same chiral transformation property, but it is straightforward to check that the combination of the first and third terms in \eqref{eq:19.5.43} behaves as a chiral-covariant derivative: that is,
\begin{equation}
\delta\mathcal D_\mu\widetilde N
=2i\boldsymbol t\cdot
[\boldsymbol{\zeta}\cdot\boldsymbol{\epsilon}]
\mathcal D_\mu\widetilde N,
\tag{19.5.46}\label{eq:19.5.46}
\end{equation}
where
\begin{equation}
\mathcal D_\mu\widetilde N
\equiv
\left[
\partial_\mu
+2i\frac{\boldsymbol t\cdot
(\boldsymbol{\zeta}\cdot\partial_\mu\boldsymbol{\zeta})}
{1+\boldsymbol{\zeta}^2}
\right]\widetilde N.
\tag{19.5.47}\label{eq:19.5.47}
\end{equation}
Thus the Lagrangian \eqref{eq:19.5.43} (and hence \eqref{eq:19.5.44}) is obviously chiral-invariant, because it is isospin-invariant, and constructed solely from the ingredients $\widetilde N$, $\mathcal D_\mu\widetilde N$, $\sigma$, and $\boldsymbol D_\mu$, all of which transform under chiral transformations with the same isospin rotation.

Again, the specific Lagrangian with which we started here is not important. As in the case of the pure pion theory considered earlier, the important thing is the chiral invariance of the Lagrangian. For chiral invariance, the Lagrangian must conserve isospin, and be constructed only from the ingredients $\widetilde N$, $\mathcal D_\mu\widetilde N$, and $\boldsymbol D_\mu$ (together with higher covariant derivatives of these objects). The most general chiral-invariant Lagrangian that is bilinear in the new nucleon field and involves no more than one derivative therefore takes the form
\begin{equation}
\mathcal L_{N,0}
=-\overline{\widetilde N}
\left[
\sl{\mathcal D}+m_N
+2ig_A\gamma_5\boldsymbol t\cdot\sl{\boldsymbol D}
\right]\widetilde N
\tag{19.5.48}\label{eq:19.5.48}
\end{equation}
or in terms of the pion field \eqref{eq:19.5.17}:
\begin{equation}
\mathcal L_{N,0}
=-\overline{\widetilde N}
\left[
\sl{\partial}+m_N
+\frac{2i\boldsymbol t\cdot
(\boldsymbol{\pi}\cdot\sl{\partial}\boldsymbol{\pi})}
{F_\pi^2[1+\boldsymbol{\pi}^2/F_\pi^2]}
+\frac{2ig_A\gamma_5\boldsymbol t\cdot
\sl{\partial}\boldsymbol{\pi}}
{F_\pi[1+\boldsymbol{\pi}^2/F_\pi^2]}
\right]\widetilde N.
\tag{19.5.49}\label{eq:19.5.49}
\end{equation}
Note that we have inserted an arbitrary constant $g_A$ in the last term in \eqref{eq:19.5.48} because this term is chiral symmetric by itself, and so chiral symmetry cannot dictate its coefficient. (This is in contrast with the third term in \eqref{eq:19.5.43}, whose coefficient is fixed by the condition that it and the first term together give a chiral invariant.) We can check that the constant $g_A$ inserted here is indeed the axial-vector coupling of beta decay by constructing the extra term in the axial-vector current that arises from the nucleonic Lagrangian \eqref{eq:19.5.48}. Alternatively, integrating by parts and using the Dirac equation, we find that the pion-nucleon interaction $-2ig_A\overline{\widetilde N}\gamma_5\boldsymbol t\cdot(\sl{\partial}\boldsymbol{\pi})\widetilde N/F_\pi$ is equivalent on the nucleon mass shell to an interaction $-4im_Ng_A\overline{\widetilde N}\gamma_5\boldsymbol t\cdot\boldsymbol{\pi}\widetilde N/F_\pi$, corresponding to a pion-nucleon coupling constant
\[
G_{\pi N}=2m_Ng_A/F_\pi,
\]
which is just the Goldberger--Treiman relation \eqref{eq:19.4.33}.

Incidentally, if we had used the Lagrangian \eqref{eq:19.5.44} in this calculation, we would have obtained the Goldberger--Treiman relation with $g_A=1$. However, this result is in fact an artifact of the particular form \eqref{eq:19.5.39} of the interaction with which we started. We might have included a non-renormalizable derivative coupling term\footnote{The left- and right-handed parts of the nucleon doublet transform according to the representations $(\tfrac12,0)$ and $(0,\tfrac12)$ of $SU(2)\cdot SU(2)$, respectively. The bilinear $\overline N\gamma_\mu\gamma_5N$ is thus a sum of terms quadratic in $(\tfrac12,0)$ or $(0,\tfrac12)$ terms, so it transforms like a direct sum of the $(1,0)$, $(0,0)$, and $(0,1)$ representations. In Eq.~\eqref{eq:19.5.50} we have coupled the $(1,0)+(0,1)$ terms in $\overline N\gamma_\mu\gamma_5N$ to the antisymmetric tensor formed from the $SU(2)\cdot SU(2)$ four-vectors $\phi_n$ and $\partial_\mu\phi_n$. Its invariance under the transformations \eqref{eq:19.5.40}, \eqref{eq:19.5.41} can of course be checked directly.}
\begin{equation}
\mathcal L'_N
=ig'\overline N
\left[
\left(
\boldsymbol t\cdot\boldsymbol{\phi}\,
\sl{\partial}\phi_4
-\phi_4\boldsymbol t\cdot
\sl{\partial}\boldsymbol{\phi}
\right)\gamma_5
+\boldsymbol t\cdot
(\boldsymbol{\phi}\cdot\sl{\partial}\boldsymbol{\phi})
\right]N.
\tag{19.5.50}\label{eq:19.5.50}
\end{equation}
It is straightforward to show that in terms of the transformed fields $\sigma$, $\widetilde N$, and $\boldsymbol{\zeta}$, this takes the form
\[
\mathcal L'_N
=-8i\sigma^2g'\overline{\widetilde N}
\boldsymbol t\cdot\sl{\boldsymbol D}\gamma_5\widetilde N.
\]
Since $\sigma$ has a non-vanishing vacuum expectation value, this makes a contribution to the pion-nucleon coupling constant and to $g_A$ proportional to $g'$. Hence the value $g_A=1$ is \emph{not} dictated by the broken $SU(2)\cdot SU(2)$ symmetry alone; by including the interaction \eqref{eq:19.5.50} and adjusting $g'$, we can make $g_A$ anything we like.

Now let us consider how to use this Lagrangian to calculate amplitudes for reactions involving both pions and nucleons. We will have to give special attention to the nucleon propagators, since a nucleon can never be a `soft' particle like a pion. A nucleon line that enters a diagram with an on-shell four-momentum $p$ of order $m_N$, and then from interactions with soft pions absorbs a net four-momentum $q$ with components much less than $m_N$, will have a propagator
\begin{equation}
\frac{-i(\sl{p}+\sl{q})+m_N}{(p+q)^2+m_N^2}
\underset{q\to0}{\longrightarrow}
\frac{-i\sl{p}+m_N}{2p\cdot q}.
\tag{19.5.51}\label{eq:19.5.51}
\end{equation}
(The neglected terms may be taken into account by including higher derivative terms in the nucleon Lagrangian.) Suppose again that all external pion four-momenta have components at most of order $Q$ and define all renormalized couplings at renormalization points of order $Q$, so that integrals converge in such a way that internal pion lines also have four-momenta $Q$. Then \eqref{eq:19.5.51} shows that internal nucleon lines make a contribution of order $1/Q$. A general Feynman diagram for such a process will make a contribution to the invariant amplitude of order $Q^\nu$, where now
\begin{equation}
\nu=\sum_iV_i(d_i+2m_i)-2I_\pi-I_N+4L.
\tag{19.5.52}\label{eq:19.5.52}
\end{equation}
Here $V_i$ is the number of vertices associated with interactions of type $i$, $d_i$ is the number of derivatives in each such interaction, $m_i$ is the number of quark mass factors in each interaction, $I_\pi$ and $I_N$ are the numbers of internal pion and nucleon lines, respectively, and $L$ is the number of loops in the graph. We use the familiar topological relations for connected graphs:
\begin{equation}
L=I_\pi+I_N-\sum_iV_i+1
\tag{19.5.53}\label{eq:19.5.53}
\end{equation}
and
\begin{equation}
2I_N+E_N=\sum_iV_i n_i,
\tag{19.5.54}\label{eq:19.5.54}
\end{equation}
where $n_i$ is the number of nucleon fields in interactions of type $i$, and $E_N$ is the number of external nucleon lines. Eliminating the quantities $I_N$ and $I_\pi$, this gives
\begin{equation}
\nu=\sum_iV_i\left(d_i+2m_i+\frac{n_i}{2}-2\right)
+2L-E_N+2.
\tag{19.5.55}\label{eq:19.5.55}
\end{equation}
The important point here is that the coefficient $d_i+2m_i+\tfrac12n_i-2$ in the first term is always positive or zero. We have already seen that $d_i+2m_i\geq2$ for the purely pionic interactions with $n_i=0$, and inspection of \eqref{eq:19.5.49} shows that $d_i\geq1$ for the pion-nucleon interactions with $n_i=2$ and $m_i=0$. Any interaction with $n_i=2$ and $m_i\geq1$ or $n_i\geq4$ clearly also has $d_i\geq2$. Hence the leading terms for $Q\ll2\pi F_\pi$ are tree graphs (that is, $L=0$) for which all interactions have
\[
d_i+2m_i+\frac{n_i}{2}-2=0.
\]

% Source: physical PDF p. 229 (printed p. 206). Complete three-panel Figure 19.3 and caption.
\begin{figure}[htbp]
  \centering
  \begin{tikzpicture}[
    x=0.82cm,
    y=0.82cm,
    line cap=butt,
    line join=round,
    nucleon/.style={line width=0.9pt},
    pion/.style={line width=0.75pt,dash pattern=on 4pt off 4pt}
  ]
    % (a): direct nucleon pole graph.
    \draw[nucleon] (0,-2.05) -- (0,2.05);
    \draw[pion] (0,0.83) -- (2.10,2.08);
    \draw[pion] (0,-0.83) -- (2.10,-2.08);
    \node at (0,-2.85) {(a)};

    % (b): crossed nucleon pole graph.
    \begin{scope}[xshift=5.25cm]
      \draw[nucleon] (0,-2.05) -- (0,2.05);
      \draw[pion] (0,0.83) -- (2.10,-1.05);
      \draw[pion] (0,-0.83) -- (2.10,1.05);
      \node at (0,-2.85) {(b)};
    \end{scope}

    % (c): pion-pion-nucleon contact graph.
    \begin{scope}[xshift=10.50cm]
      \draw[nucleon] (0,-2.05) -- (0,2.05);
      \draw[pion] (0,0) -- (2.10,1.22);
      \draw[pion] (0,0) -- (2.10,-1.22);
      \node at (0,-2.85) {(c)};
    \end{scope}
  \end{tikzpicture}
  \renewcommand{\thefigure}{19.3}
  \caption{Feynman diagrams used with an effective chiral Lagrangian to calculate the scattering of a soft pion from a nucleon. Dashed lines are pions; solid lines are nucleons.}
  \label{fig:19.3}
\end{figure}


The interactions that satisfy this condition are just those shown explicitly in Eqs. \eqref{eq:19.5.23}, \eqref{eq:19.5.34}, and \eqref{eq:19.5.49}, plus possible interactions with $d_i=0$ and $n_i=4$:
\begin{equation}
\left(\overline{\widetilde N}\Gamma_\alpha\widetilde N\right)
\left(\overline{\widetilde N}\Gamma^\alpha\widetilde N\right),
\tag{19.5.56}\label{eq:19.5.56}
\end{equation}
where $\Gamma_\alpha$ and $\Gamma^\alpha$ are any matrices in spin and isospin space that yield Lorentz, space inversion, and isospin-invariant four-fermion interactions. These last interactions are important for multinucleon processes~\hyperref[ch19-ref-29]{[29]}, which will not be considered here.

Let us now apply this method to pion-nucleon scattering. For pion energies roughly of order $m_\pi$, the pion-nucleon scattering amplitude is given by the Feynman graphs of Figure~\ref{fig:19.3}, all of which make contributions of order $m_\pi$ to the invariant amplitude. However, at threshold the leading contribution is from Figure~\ref{fig:19.3}(c), the others being suppressed by an extra factor of $m_\pi/m_N$. This is because at threshold in the rest frame the incoming and outgoing pion four-momenta $q,q'$ and nucleon four-momenta $p,p'$ are given by
\begin{equation}
q=q'=(m_\pi,\boldsymbol{0})
=\left(\frac{m_\pi}{m_N}\right)p
=\left(\frac{m_\pi}{m_N}\right)p'.
\tag{19.5.57}\label{eq:19.5.57}
\end{equation}

Thus the invariant amplitude at threshold from either Figure~\ref{fig:19.3}(a) or \ref{fig:19.3}(b) is proportional to
\begin{align*}
&\overline u\gamma_5\sl{q}\,
\frac{-i(\sl{p}\mathbin{\pm}\sl{q})+m_N}
{(p\mathbin{\pm}q)^2+m_N^2}
\gamma_5\sl{q}\,u
\\
&\quad=
\frac{m_\pi^2\overline u\gamma_5
\{-i\sl{p}(1\mathbin{\pm}m_\pi/m_N)+m_N\}
\gamma_5u}
{\mp2m_\pi m_N-m_\pi^2}
\\
&\quad=
\frac{\mp m_\pi}{2m_N\mathbin{\pm}m_\pi}
\overline u\{i\sl{p}(1\mathbin{\pm}m_\pi/m_N)+m_N\}u
=\frac{m_\pi^2}{2m_N\mathbin{\pm}m_\pi},
\end{align*}
where $u$ is a Dirac spinor with $\overline uu=1$. In contrast, the diagram of Figure~\ref{fig:19.3}(c) gives a contribution to $M$ of order $m_\pi$. We may write this amplitude as a $2\cdot2$ matrix in the nucleon isospin indices:
\[
M_{ba}
=-\frac{2i}{F_\pi^2}t_c\epsilon_{abc}
\overline u(-i\sl{q}-i\sl{q}')u,
\]
where $q'$ and $q$ are the final and initial pion four-momenta and $b$ and $a$ are the corresponding isovector indices. Using \eqref{eq:19.5.57} and the momentum-space Dirac equation $(i\sl{p}+m_N)u=0$, this becomes
\begin{equation}
M_{ba}
=-\frac{4im_\pi}{F_\pi^2}t_c\epsilon_{abc}
=-\frac{4m_\pi}{F_\pi^2}
\boldsymbol t\cdot[\boldsymbol t^{(\pi)}]_{ba},
\tag{19.5.58}\label{eq:19.5.58}
\end{equation}
where $[t_c^{(\pi)}]_{ba}\equiv-i\epsilon_{bac}$ is the pion isovector matrix. The matrix $\boldsymbol t\cdot\boldsymbol t^{(\pi)}$ has eigenvalues in states of total isospin $T$ equal to $\tfrac12[T(T+1)-2-\tfrac34]$, so in the two isospin states, $T=1/2$ and $T=3/2$, the invariant amplitudes are~\hyperref[ch19-ref-30]{[30]}
\[
M_{1/2}=4m_\pi/F_\pi^2,
\qquad
M_{3/2}=-2m_\pi/F_\pi^2.
\]
These results are usually expressed in terms of the scattering lengths, which are defined as the invariant amplitudes divided by $4\pi(1+m_\pi/m_N)$:
\begin{equation}
a_{1/2}
=\frac{m_\pi}{\pi F_\pi^2(1+m_\pi/m_N)}
=0.15m_\pi^{-1},
\tag{19.5.59}\label{eq:19.5.59}
\end{equation}
\begin{equation}
a_{3/2}
=-\frac{m_\pi}{2\pi F_\pi^2(1+m_\pi/m_N)}
=-0.075m_\pi^{-1}.
\tag{19.5.60}\label{eq:19.5.60}
\end{equation}
These are in reasonable agreement with the experimental values~\hyperref[ch19-ref-12]{[12]} $a_{1/2}=(0.173\mathbin{\pm}0.003)m_\pi^{-1}$ and $a_{3/2}=(-0.101\mathbin{\pm}0.004)m_\pi^{-1}$. (The results \eqref{eq:19.5.59} and \eqref{eq:19.5.60} are only supposed to be valid to lowest order in $m_\pi/m_N$, but we retain the factor $1+m_\pi/m_N$, because it arises just from the definition of the scattering lengths.)

The corrections to the scattering lengths of next order in $m_\pi$ arise from several sources. There are the Born diagram graphs of Figures~\ref{fig:19.3}(a) and \ref{fig:19.3}(b), which are nominally of leading order but as we have seen at threshold are suppressed by an extra factor of $m_\pi/m_N$. There are additional tree graphs containing a vertex with two derivatives. Of special interest are tree graphs containing a vertex with no derivatives that arises from a symmetry-breaking interaction proportional to $m_\pi^2$. These interactions must have the chiral and spacetime transformation properties of the operators in \eqref{eq:19.4.34}: the fourth and third components of two different chiral four-vectors $\Phi_n^+$ and $\Phi_n^-$. There are two obvious candidates for such operators that are bilinear in the nucleon fields: a $\Phi_4^+$ term,
\[
\overline NN
=\left(\frac{1-\boldsymbol{\zeta}^2}
{1+\boldsymbol{\zeta}^2}\right)
\overline{\widetilde N}\widetilde N
-4i\left(\frac{\boldsymbol{\zeta}}
{1+\boldsymbol{\zeta}^2}\right)
\cdot\overline{\widetilde N}\gamma_5\boldsymbol t\widetilde N,
\]
and a $\Phi_3^-$ term,
\[
\overline Nt_3N
=\overline{\widetilde N}t_3\widetilde N
-2\left(\frac{\zeta_3}{1+\boldsymbol{\zeta}^2}\right)
\overline{\widetilde N}\boldsymbol t\cdot
\boldsymbol{\zeta}\widetilde N
-i\left(\frac{\zeta_3}{1+\boldsymbol{\zeta}^2}\right)
\overline{\widetilde N}\gamma_5\widetilde N.
\]
Chiral symmetries act on $\widetilde N$ like ordinary isospin rotations, so they cannot mix up scalar and pseudoscalar nucleon bilinears. Thus there are really two independent $\Phi_4^+$ operators:
\begin{equation}
\left(\frac{1-\boldsymbol{\zeta}^2}
{1+\boldsymbol{\zeta}^2}\right)
\overline{\widetilde N}\widetilde N
\tag{19.5.61}\label{eq:19.5.61}
\end{equation}
and
\begin{equation}
i\left(\frac{\boldsymbol{\zeta}}
{1+\boldsymbol{\zeta}^2}\right)
\cdot\overline{\widetilde N}\gamma_5\boldsymbol t\widetilde N,
\tag{19.5.62}\label{eq:19.5.62}
\end{equation}
and two independent $\Phi_3^-$ operators:
\begin{equation}
\overline{\widetilde N}t_3\widetilde N
-2\left(\frac{\zeta_3}{1+\boldsymbol{\zeta}^2}\right)
\overline{\widetilde N}\boldsymbol t\cdot
\boldsymbol{\zeta}\widetilde N
\tag{19.5.63}\label{eq:19.5.63}
\end{equation}
and
\begin{equation}
i\left(\frac{\zeta_3}{1+\boldsymbol{\zeta}^2}\right)
\overline{\widetilde N}\gamma_5\widetilde N.
\tag{19.5.64}\label{eq:19.5.64}
\end{equation}
We shall show in the next section that these are the only operators with the transformation properties of the terms in \eqref{eq:19.4.34} that are bilinear in the nucleon fields. The operators \eqref{eq:19.5.62} and \eqref{eq:19.5.64} evidently provide isospin-conserving and isospin-violating corrections to the Goldberger--Treiman formula for the pion-nucleon couplings. The other two operators, \eqref{eq:19.5.61} and \eqref{eq:19.5.63}, contribute directly to both the nucleon mass and to low energy pion-nucleon scattering. From their contribution to the nucleon mass, we see that these latter terms enter into the effective Lagrangian in the form (now replacing $\boldsymbol{\zeta}$ with the conventionally normalized pion field):
\begin{align}
\delta\mathcal L_{\mathrm{eff}}
={}&-\frac{\delta m_p+\delta m_n}{2}
\left(\frac{1-\boldsymbol{\pi}^2/F_\pi^2}
{1+\boldsymbol{\pi}^2/F_\pi^2}\right)
\overline{\widetilde N}\widetilde N
\notag\\
&-(\delta m_p-\delta m_n)
\left[
\overline{\widetilde N}t_3\widetilde N
-\frac{2}{F_\pi^2}
\left(\frac{\pi_3}{1+\boldsymbol{\pi}^2/F_\pi^2}\right)
\overline{\widetilde N}\boldsymbol t\cdot
\boldsymbol{\pi}\widetilde N
\right],
\tag{19.5.65}\label{eq:19.5.65}
\end{align}
where $\delta m_p$ and $\delta m_n$ are the contributions of the quark mass terms \eqref{eq:19.4.34} to the proton and neutron masses. This makes a contribution to the pion-nucleon scattering amplitude (again written as a matrix in the isospin space of the nucleon):
\begin{equation}
\delta M_{ba}
=\frac{2[\delta m_p+\delta m_n]}{F_\pi^2}\delta_{ab}
+\frac{2[\delta m_p-\delta m_n]}{F_\pi^2}
(t_a\delta_{3b}+t_b\delta_{3a}).
\tag{19.5.66}\label{eq:19.5.66}
\end{equation}
The first term in \eqref{eq:19.5.65} is often known as the `$\sigma$-term'. The second term is an isospin violating correction to the $\sigma$-term, which shows up only in processes involving neutral pions, such as the charge-exchange processes $\pi^++n\to\pi^0+p$ and $\pi^-+p\to\pi^0+n$. Even though it is not possible to measure the nucleon mass shifts $\delta m_{p,n}$ directly, we shall see in Section 19.7 that an $SU(3)$ symmetry allows their difference to be calculated from hyperon mass differences. This yields the result $\delta m_p-\delta m_n=-2.5\ \mathrm{MeV}$. Unfortunately, we still do not have a firm theoretical estimate of the coefficient $\delta m_p+\delta m_n$ of the first term in \eqref{eq:19.5.65}.

Eq.~\eqref{eq:19.4.34} shows that in quantum chromodynamics the coefficients $\delta m_p+\delta m_n$ and $\delta m_p-\delta m_n$ in \eqref{eq:19.5.66} are respectively proportional to $m_u+m_d$ and $m_u-m_d$. We shall see in Section 19.7 that $m_u$ and $m_d$ are not at all degenerate, so these coefficients are roughly of the same order of magnitude; the isospin-violating corrections to the $\sigma$-term are not much smaller than the $\sigma$-term itself. We see again that the reason that isospin conservation is such a good approximation in hadronic physics is not that the $u$ and $d$ quark masses are nearly equal, but just that they are small. The use here of Lagrangians like \eqref{eq:19.5.19} or \eqref{eq:19.5.49} is one example of the method of \emph{effective field theories}, already introduced in Section 12.3. Similar techniques have been employed to deal with meson and baryon interactions including strange particles (see Section 19.7), with quark and lepton interactions at energies below the scale of electroweak symmetry breaking (see Section 21.3), and even with superconductivity (see Section 21.6). In all these cases effective field theories provide the most convenient method for working out the consequences of symmetries and the general principles underlying quantum field theory.

\begin{center}
$\ast\quad\ast\quad\ast$
\end{center}

The low energy theorems provided by broken symmetries when combined with dispersion relations yield useful sum rules. Let's see how this works for chiral symmetry, ignoring the small $u$ and $d$ quark masses and the pion mass. Consider the forward scattering of a massless pion with isovector index $a$ and four-momentum $q$ on a nucleon of four-momentum $p$, yielding a pion with isovector index $b$. The scattering amplitude $M$, defined by $S_{fi}=-2\pi i\delta^4(P_f-P_i)M_{fi}$, is given to first order in $q$ by the Feynman diagrams of Figure~\ref{fig:19.3} as
\begin{align}
-2\pi iM
={}&\frac{1}{(2\pi)^6(2q^0)}
\overline u\Bigg[
\notag\\
&\qquad
\left(-i\frac{(2\pi)^4\,2g_A\gamma_5t_b\sl{q}}{F_\pi}\right)
\left(\frac{-i}{(2\pi)^4}
\frac{-i(\sl{p}+\sl{q})+m_N}{(p+q)^2+M^2}\right)
\left(i\frac{(2\pi)^4\,2g_A\gamma_5t_a\sl{q}}{F_\pi}\right)
\notag\\
&\qquad+
\left(i\frac{(2\pi)^4\,2g_A\gamma_5t_a\sl{q}}{F_\pi}\right)
\left(\frac{-i}{(2\pi)^4}
\frac{-i(\sl{p}-\sl{q})+m_N}{(p-q)^2+M^2}\right)
\left(-i\frac{(2\pi)^4\,2g_A\gamma_5t_b\sl{q}}{F_\pi}\right)
\notag\\
&\qquad-\frac{4it_c\epsilon_{abc}\sl{q}}{F_\pi^2}
\Bigg]u.
\tag{19.5.67}\label{eq:19.5.67}
\end{align}
To this order in $q$, the nucleon propagator in the first two terms inside the square brackets may by approximated by
\[
\frac{-i(\sl{p}\mathbin{\pm}\sl{q})+m_N}
{(p\mathbin{\pm}q)^2+M^2}
\approx
\frac{-i\sl{p}+m_N}{\mathbin{\pm}2p\cdot q}.
\]
Eq.~\eqref{eq:19.5.67} can then be further simplified, using the relations $\sl{q}\sl{q}=q^2=0$, $\overline u\gamma^\mu u=-ip^\mu/m_N$, and $[t_a,t_b]=i\epsilon_{abc}t_c$. Also, in the laboratory frame (which for small $q$ is the same as the center-of-mass frame) we have $p\cdot q=-m_N\omega$, where $\omega\equiv q^0$ is the pion energy in the nucleon rest frame. Finally, the conventional forward scattering amplitude (whose absolute square is the differential cross section for forward scattering) is given by Eq.~(3.6.9) as $f=-4\pi^2\omega M$. Putting this all together, we find the forward scattering amplitude at pion energy $\omega$ is given for $\omega\ll m_N$ by
\begin{equation}
f_{ba}(\omega)
\longrightarrow
-i\frac{\omega}{\pi F_\pi^2}(1-g_A^2)\epsilon_{abc}t_c.
\tag{19.5.68}\label{eq:19.5.68}
\end{equation}
In particular, for $\pi^+$--proton scattering, we must contract \eqref{eq:19.5.68} with $v_b^*v_a$, where $\boldsymbol v$ is the normalized isovector $(1,i,0)/\sqrt{2}$, and set $t_3=+\tfrac12$. The forward low energy $\pi^+$--proton scattering amplitude is then
\begin{equation}
f_{\pi^+p}(\omega)
\longrightarrow
-\frac{\omega}{2\pi F_\pi^2}(1-g_A^2).
\tag{19.5.69}\label{eq:19.5.69}
\end{equation}

Now, the dispersion relation for the forward scattering of a massless $\pi^+$ on a proton is given by Eq.~(10.8.24) (with subtraction polynomial $P(E)\propto E$) as
\begin{align}
f_{\pi^+p}(\omega)
={}&R+\frac{i\omega}{4\pi}\sigma_{\pi^+p}(\omega)
\notag\\
&+\frac{\omega}{4\pi^2}\int_0^\infty
\left[
\frac{\sigma_{\pi^+p}(E)}{E-\omega}
-\frac{\sigma_{\pi^-p}(E)}{E+\omega}
\right]dE,
\tag{19.5.70}\label{eq:19.5.70}
\end{align}
where $R$ is a possible subtraction constant, and $\omega$ and $E$ are pion energies in the proton rest frame. Comparison of this dispersion relation with the low energy limit \eqref{eq:19.5.69} shows that $R=0$, and that
\begin{equation}
g_A^2
=1+\frac{F_\pi^2}{2\pi}\int_0^\infty
[\sigma_{\pi^+p}(E)-\sigma_{\pi^-p}(E)]\,dE/E.
\tag{19.5.71}\label{eq:19.5.71}
\end{equation}
This is the celebrated Adler--Weisberger sum rule~\hyperref[ch19-ref-19]{[19]}, for the case of exact chiral symmetry, with $m_\pi=0$. A similar sum rule may be derived for the scattering of pions on any baryon or meson, including the pion itself~\hyperref[ch19-ref-30a]{[30a]}.
